\section{Discussion}\label{sec:disc}
% =====================================================================
% Section 8: Discussion.
% Sources (final records only): research/00-brief.md (questions, scope); the section files of this paper
% (every claim below points to a labelled result there); research/tracks/*/notes-final.md (open problems:
% single \S13, cases \S\S5--6, untagged \S10, experiments \S\S5--6); research/prior/induction/second-order-notes.md
% \S3 (literature: Vaught, metarules, Libal--Miller, proof assistants) and raw-impossibility-notes.md B7
% (Ryll-Nardzewski, Vaught). New bibliography entries are in bib/disc.bib, each confirmed by web search.
% Review of the paper draft (research/paper-review/, edits/E1-results.txt, E2-results.txt): the claims below were
% re-read against the edited sections; in particular the omega-step is "in general not derivable", separation of PA,
% ZF and ZFC is qualified as in many.tex, R-Delta_0 decides truth only for quantifier-free phi and at unbounded depth,
% the sharing-pass problem follows prop:exp:share as corrected, and the SO problem cites conj:zf:root.
% No new results are claimed here.
% =====================================================================

\Cref{sec:univ,sec:zf,sec:many} answer the three questions and \cref{sec:exp} tests the answers. This section says
what they mean for the project they belong to, learning which inference rules a community uses from imitation,
coherence and the world \citep{claude2026whatfollows}. It then relates them to earlier work, collects the open
problems, and states limitations and scope. It claims no new results.

% ---------------------------------------------------------------------
\subsection{What the answers mean}\label{sec:disc:meaning}

\paragraph{An instance schema is not a universal sentence.}
Instances $\varphi(t)$ at closed terms determine the instance schema $\varphi(z)$ (for one variable, two instances whose
terms have different head symbols suffice; \cref{thm:univ:anchor}), but not the sentence $\forall x\varphi$. The step to
$\forall x\varphi$ is an $\omega$-rule step: truth-safe for every $\varphi$ iff the closed-term substructure is
elementary (\cref{prop:univ:tv}), and in general not derivable, even over $\Q$ (\cref{ex:univ:q}). A cautious learner
whose class admits parameters as values takes it silently after the first anchor (\cref{prop:univ:open}), and closed
data cannot show whether a community takes it. A closedness guard blocks the step; data with a parameter at every
variable license it as ordinary generalization (\cref{tab:univ:learner}). Without labels, $\forall x\varphi$ is an
unrefuted merge of instances (\cref{prop:many:forall}), certified by the absence of a counterexample, not by a
derivation. With $\Delta_0$ evaluation ``never refuted'' means ``true'' for quantifier-free $\varphi$, though only in the
limit of unbounded depth, and not for every quantified $\varphi$ (\cref{prop:many:forall,prop:many:qf}).

\paragraph{First-order learnability depends on the formulation.}
A ZF schema is a guarded first-order pattern iff its metavariable is applied to literally the same argument tuple at
every occurrence in the chosen syntax (\cref{thm:zf:classify,prop:zf:alpha}). So $\in$-induction passes with de Bruijn
indices and fails with names, Collection the other way round, and $\alpha$-renaming the frame changes the verdict
(\cref{tab:zf:classify}). No finite union of first-order schemas covers raw PA induction truth-soundly
(\cref{thm:zf:indunion}), yet the form $\IndEq$ is a single guarded first-order pattern (\cref{prop:zf:indeq}). A
checker should learn in a class whose verdicts do not depend on what a community regards as notation. The $\lambda$
convention does: bound-variable names do not matter, freshness is built in, and capture cannot occur. In a first-order
syntax capture can make an unguarded learner unsound outright, and a guard must exclude it exactly when capture
instances exist (\cref{prop:univ:capture}, \cref{lem:univ:capcrit}).

\paragraph{Why the pattern fragment matters.}
Every ZF schema, in every formulation considered, is a higher-order pattern \citep{miller1991logic}: $L_\in$ has no
function symbols, so an instance only renames variables inside the body. Here existing technology suffices: the pattern
lgg recovers each schema from any anchor (\cref{cor:zf:patlgg}), and for one metavariable the anchors, root variation
plus non-vacuity, are the same in every class between $\PAT$ and $\SO$ (\cref{thm:zf:anchor}). With several
metavariables, as in untagged mixtures, a Plotkin-(D)-type event appears and differs between $\PAT$, $\DT$ and $\SO$
(\cref{thm:zf:multi}): whether data are an anchor depends on the class.

\paragraph{Induction needs derived occurrences.}
In $T_{\Ind}$ the occurrences $P(0)$ and $P(Sx)$ apply the motive to terms, so their values are \emph{derived} from the
value at the pattern occurrence $P(x)$. First-order and pattern anti-unification are both unsound on induction
(\cref{prop:zf:indtwo,prop:zf:indpat}), and some unsound first-order lggs are never refuted by true closed
quantifier-free sentences (\cref{thm:zf:K0}). $\DT$ is built for this case: a pattern occurrence for every metavariable
and metavariable-free arguments keep matching unique and linear and the minimal covering templates finitely many
(\cref{thm:setting:matching,thm:single:min}). Without determinacy anchors must be stronger (\cref{prop:zf:indso}); with
metavariables inside arguments minimal templates can be infinitely many (\cref{rem:setting:nested}). The price is
non-unitarity, up to $4^n$ minimal templates (\cref{prop:single:blowup}), and caution pays it: the verifier intersects
them all by a polynomial feature test (\cref{thm:single:feat}), which on the data behind $K_0$ rejects its false
instance (\cref{sec:many:examples}).

\paragraph{Refutation is the source of labels.}
Without a bound on the number of schemas the
cautious learner accepts only its data (\cref{prop:many:bound}); with a bound, spare slots cost data diversity that no
sound negatives replace (\cref{thm:many:splits}). Refutations remove over-general \emph{lumps}, a bound removes
\emph{splits}, which are sound and so invisible to refutation. \DTRC{} uses no bound. Under refutation separation it
recovers the hidden labels and is then the tagged learner, with tagged anchors and rates (\cref{thm:many:dtrc}). So the
world and coherence channels, which in the parent report rule out errors, here also do the work of citations. They
separate $\Q$'s axioms and $T_{\Ind}$ at small depth (proved), and ZF and ZFC on every sampled cross pair and every run
(computed; \cref{sec:many:examples}); for the mixtures of \cref{sec:exp} separation is proved for an ideal refuter
(\cref{cor:exp:mixsep}). The depth cannot be eliminated (\cref{thm:many:depth}), and compression does not replace
refutation: MDL \citep{rissanen1978modeling} and the size principle \citep{tenenbaum2001generalization} follow the
statistics of usage, not the logical boundaries of schemas (\cref{sec:many:mdl}).

\paragraph{Truth-soundness and target-soundness come apart.}
``Which axioms does this community use?'' asks for target-soundness, ``is it safe to accept this step?'' for
truth-soundness, and refutation serves only the second. Merges without false instances are never refuted
(\cref{lem:setting:onesided}): $\forall x\varphi$ from its instances, the true template $0+z=z$ from data of two targets
(\cref{thm:many:residue}), the valid residual templates of the experiments (\cref{lem:exp:valid}). They can leave
\DTRC{} truth-sound but not target-sound, unavoidably (\cref{prop:many:ambiguity}); target-soundness needs labels, a known
bound (\cref{rem:many:combine}) or the community's own judgments. Truth-soundness in turn holds only up to a residue of
unrefuted unsound merges, which has natural members: the universal-set schema from weak Union and Power, when only logic
and hereditarily finite counterexamples are used (\cref{prop:many:univset}); $K_0$-shaped templates, once a mistake
removes the sound alternative (\cref{ex:many:noise}); and, at small refutation budgets, a false sentence accepted from
data with mistakes (\cref{prop:exp:qF}). This is the parent report's picture in a sharp case: every signal is
one-sided, and what can be proved is soundness with a declared residue.

% ---------------------------------------------------------------------
\subsection{Relation to the literature}\label{sec:disc:lit}

\paragraph{Anti-unification and inductive logic programming.}
The first-order core is the least general generalization of \citet{plotkin1970note} and
\citet{reynolds1970transformational}, the basis of bottom-up generalization in inductive logic programming
\citep{muggleton1991inductive}. Plotkin's witness events are our anchor conditions for first-order targets, in $\DT$
too (\cref{cor:single:special}), and the feature test (\cref{thm:single:feat}) is the second-order form of his rule
that equal columns get equal variables. Metarules in meta-interpretive learning are second-order Horn clauses given in
advance \citep{muggleton2015meta,cropper2020logical}; here the second-order templates are what is learned.

\paragraph{Higher-order anti-unification and matching.}
For higher-order patterns the lgg is unique \citep{pfenning1991unification} and computable in linear time
\citep{baumgartner2017higher}. Second-order generalization beyond patterns is studied by
\citet{hirata2004generalization} and surveyed by \citet{cerna2023antiunification}; we have not checked whether $\DT$
or the feature test appears there. $\DT$ matching is Miller's pattern unification with one side ground, followed by
evaluation (\cref{thm:setting:matching}; general second-order matching: \cref{app:setting:matching}).
The closest relatives we know of are the deterministic second-order patterns of \citet{yokoyama2004deterministic},
which relax Miller's restriction while keeping matching deterministic, and the functions-as-constructors fragment of
\citet{libal2016functions}; \citet{niederhauser2026unification} give a unification procedure for the former that
generalizes the latter. We know these classes only from abstracts and secondary descriptions, by which, as we read them,
every argument of a metavariable must contain a bound variable, which the argument of $P(0)$ does not. We have compared neither class
with $\DT$.

\paragraph{Learning from positive data.}
By a theorem of \citet{gold1967language}, an untagged learner needs a bound or a restriction on the practice such as
refutation separation (\cref{prop:many:bound}). Every anchor is a tell-tale in the sense of
\citet{angluin1980inductive}, and up to effectivity anchors are the $\subseteq$-tell-tales of strong-monotonic
learning \citep{lange1992types}; the cautious verifier is reliable learning \citep{rivest1988learning}, or perfect
selective classification \citep{elyaniv2010foundations}, applied to sentences. First-order patterns are the tree
analogue of Angluin's pattern languages \citep{angluin1980finding}. Finite elasticity and its preservation under bounded
unions \citep{wright1989identification,motoki1991correct} give anchors for unions of determinate templates
(\cref{cor:single:unions}). \Citet{shinohara1991inductive} shows learnability from positive data for classes defined by
length-bounded elementary formal systems, and \citet{arimura1994finding} compute minimal multiple generalizations for
unions of pattern languages; we have not compared their complexity results with \cref{thm:many:complexity}. What the
present setting adds is soundness at all times against adaptive provers, explicit anchors with rates, and refutation
as negative evidence.

\paragraph{Axiomatizability by schemas.}
\citet{vaught1967axiomatizability} showed that every recursively enumerable theory that directly interprets a weak set
theory is axiomatizable by a schema, a formula with a new relation symbol whose substitution instances are the axioms.
Such a schema is a second-order template, with induction as the paradigm; $\DT$ is a version with unique matching,
learnable from instances. That PA is not finitely axiomatizable \citep{rylln1952axiomatizability} concerns
consequence, whereas \cref{thm:zf:indunion} concerns the syntax of substitution; neither implies the other, since PA is
$\Q$ plus one guarded first-order pattern (\cref{prop:zf:indeq}).

\paragraph{Proof assistants.}
Formal libraries record which axiom, lemma or recursor a step applies, as the tagged setting assumes, and they record or
recompute motives. Lean~4's \texttt{induction} tactic records the motive, an implicit argument of the recursor, as a
$\lambda$ in the kernel term \citep{demoura2021lean4}: the $\DT$ matcher. Metamath's set.mm states induction with
implicit-substitution hypotheses and distinct-variable conditions \citep{megill2019metamath}: a first-order pattern with
freshness guards that are sound but not minimal (\cref{sec:zf:pa}, \cref{app:zf:pa}). So for formal mathematics the
tagged results are the relevant ones; the untagged theory matters where citations are missing, as in informal
mathematics.

% ---------------------------------------------------------------------
\subsection{Open problems}\label{sec:disc:open}

The research records leave the following problems open; gray tags name the record, pointers the nearest result. Two
problems listed as open in one record are settled in another: finiteness of minimal covering templates in $\DT$
(\cref{thm:single:min}) and the complexity of the refutation test (\cref{thm:many:complexity}; but see
\ref{op:disc:cost}).
\begin{enumerate}[label=O\arabic*,ref=O\arabic*,leftmargin=*,itemsep=0pt,topsep=3pt]
\item The anchor condition without \Rich{} (\cref{thm:single:anchor}). Does $\inst(T)\supseteq\inst(T')$ imply
$T\gen T'$ in $\DT$? (The verifier does not need it.) \src{single \S13}
\item A theory of parameters invariant under renaming: \cref{prop:setting:align} covers only targets whose parameters
all occur in their skeleton (\cref{rem:setting:alignfail}) \src{single \S13; paper review}.
\item Is the number of equation kills linear (\cref{conj:single:kill})? Escalations are $\Theta(N^2)$ in general
(\cref{prop:single:quadratic}); the counts for the ZF schemas are not known \src{single \S13; cases \S6}.
\item Size-bounded subclasses of $\DT$, in which induction is closed iff $s\ge12$ (\cref{prop:zf:indgen})
\src{prior induction C6.2}; enumeration of $\Min(D)$ with polynomial delay, and counting it
(\cref{prop:single:blowup}); lower bounds for the coincidence terms of the rates with infinite support
(\cref{prop:single:kappa}) \src{single \S13}.
\item Anchors in $\SO$: is root-level distinctness necessary with several metavariables (conjectured not,
\cref{conj:zf:root})? What is the exact condition for instance schemas over languages with closed terms
(\cref{prop:univ:dt}(c))? \src{cases \S6}
\item How $\DT$ relates to deterministic second-order patterns and to functions-as-constructors
(\cref{sec:disc:lit}) \src{prior induction \S3}.
\item Does ZF without Foundation prove Collection? It decides whether the named lgg of anchor data of spelled-out
Replacement is truth-sound there (\cref{prop:zf:coll}); unknown to us, it may be known. The lgg of non-anchor $\ReplS$
data when guards are also learned for term metavariables (\cref{app:zf:fofail}). Partial or inconsistent
$\alpha$-renaming of binders (\cref{prop:zf:alpha} covers full renaming) \src{cases \S\S5--6}.
\item Spare-slot thresholds beyond one-variable schemas, e.g.\ for $T_{\Ind}$ with $m\ge9$ slots
(\cref{rem:many:numerals}); usable bounds on the separation depth for natural targets, where no uniform computable
bound exists (\cref{thm:many:depth}) \src{untagged \S10}.
\item A natural permanent residue, i.e.\ an unsound, never-refuted merge of natural axioms; in arithmetic it needs
quantifiers in the template (\cref{prop:many:qf}) \src{untagged \S10}.
\item A merge rule that provably prevents fragmentation, which needs more than anchors
(\cref{prop:many:ambiguity,prop:many:noise}); guarantees for bootstrapped coherence when only some clusters are pure
(\cref{prop:many:univset}); a natural universal code that provably selects the target partition under a well-specified
usage model (\cref{sec:many:mdl}) \src{untagged \S10}.
\item Refutation separation for similar targets, characterized by the targets alone
(\cref{prop:exp:hard,lem:exp:valid}); a sharing pass with guarantees for every target (\cref{prop:exp:share} covers only
targets that label a phase-one cluster, \cref{ex:exp:share}, and needs (H4$^+$) and (H5)) \src{experiments \S6}.
\item\label{op:disc:cost} Is the refutation test NP-hard for true data and false negatives, as in \DTRC{} (the proof of
\cref{thm:many:complexity} uses a false datum)? And the union verifier for bounded binder depth and fixed $k\ge3$?
\src{single \S13}
\item A natural refuter that refutes every false minimal template of a cross pair within a budget polynomial in the
data (\cref{prop:exp:qF}); oracles that construct witnesses (\cref{sec:exp:impl}) \src{experiments \S6}.
\end{enumerate}
Beyond the records, the limitations below suggest typed object languages (the motive of Lean's recursor is
dependently typed), inference rules with premises and contexts, which enter here only through encodings such as Sub,
and real corpora such as set.mm or a Lean library as data.

% ---------------------------------------------------------------------
\subsection{Limitations}\label{sec:disc:limits}

We read Question~1's ``sentences of the form $\varphi(x)$'' as instances at terms; read literally, such a datum is
already the axiom (\cref{sec:univ}). Object languages are single-sorted and first-order, data are in closure-normal form,
and the theory treats parameters as constants while the implementation aligns them (\cref{sec:setting:params}). The
practice is assumed correct outside the noise analysis (\cref{prop:many:noise}), and rates assume i.i.d.\ data, while
real usage is heterogeneous (\cref{sec:many:mdl}). The implemented class $\DTF$ lacks term metavariables of positive
arity, so it can accept more than $\DT$ (E6), and the feature test is not implemented; the oracles are incomplete, the
refuter is budgeted and history-dependent, and the data are synthetic and small (\cref{sec:exp:limits}). Each research
record was checked by an adversarial referee with independent code. A review of the paper draft found two more false
propositions, now restricted (\cref{prop:setting:align,prop:exp:share}), and several overclaims (\cref{app:ver:corr}).
Results added after the referee round, including the material written during that review, were checked by proof and
computation but not attacked by an independent referee (\cref{app:ver:unref}). Bounded searches are evidence, not
proof; textbook wordings of the ZF axioms are recalled, not checked against the books; and overlaps with the literature
marked above as unchecked remain so.

% ---------------------------------------------------------------------
\subsection{Scope: checking, not proving}\label{sec:disc:scope}

Everything here concerns checking: deciding whether a proposed sentence is licensed by the axioms and schemas a
community has been seen to use. The learner never searches for proofs, and the prover of the protocol is arbitrary and
may be adversarial (\cref{sec:setting:learning}). The guarantees are soundness against every prover at all times,
relative to the class and, without labels, to a residue, and exactness after an anchor, with rates; the costs are data
and verification costs, not search costs. Nothing here makes a prover stronger, and an accepted sentence is one the
community's practice licenses, which can differ from what is true. Possible uses, not tested here, are auditing which
axioms a body of uncited proofs relies on (for instance Collection rather than Replacement, or Choice) and sending to
human review the steps that no learned schema licenses, the escalations of the protocol.
